\setchoices{label=(\Alph*)}

\begin{question}[index=1]
设函数 $f(x)=\dfrac{x}{\tan x}\ln\left|x-\dfrac{\pi}{2}\right|$，则 $f(x)$ 在 $(-\pi,\pi)$ 内有
\begin{choices}[columns=1]
\item 一个可去间断点，两个跳跃间断点
\item 两个可去间断点，一个跳跃间断点
\item 三个可去间断点
\item 两个无穷间断点，一个可去间断点
\end{choices}
\end{question}
\QuestionBreak

\begin{question}[index=2]
设 $u_0=0,\ u_{n+1}=\sqrt{2+u_n}$，则级数 $\displaystyle\sum_{n=1}^{\infty}(-1)^n\sqrt{2-u_n}$
\begin{choices}[columns=4]
\item 条件收敛
\item 绝对收敛
\item 发散
\item 无法判断
\end{choices}
\end{question}
\QuestionBreak

\begin{question}[index=3]
设 $\displaystyle I_k=\iint_{D_k}\bigl[xy^2\eu^{-x^2}+x^2\sin(x^2+y^2)\bigr]\dd\sigma$，$k=1,2,3$，其中
$D_1:x^2+y^2\leq\dfrac{\pi}{2}$，$D_2:x^2+y^2\leq\pi$，$D_3:x^2+y^2\leq2\pi$，则 $I_1,I_2,I_3$ 三者的大小关系为
\begin{choices}[columns=4]
\item $I_1\leq I_2\leq I_3$
\item $I_3\leq I_2\leq I_1$
\item $I_1\leq I_3\leq I_2$
\item $I_3\leq I_1\leq I_2$
\end{choices}
\end{question}
\QuestionBreak

\begin{question}[index=4]
设 $y=f(x)$ 是微分方程 $y'-P(x)y=\dfrac12\eu^x-1$ 满足 $y(0)=1$ 的解，其中 \allowbreak\mbox{$\displaystyle P(x)=\int_0^1|x-t|\dd t$，则有}
\begin{choices}[columns=1]
\item 在点 $x=0$ 处 $f(x)$ 取极大值
\item 在点 $x=0$ 处 $f(x)$ 取极小值
\item $(0,f(0))$ 为曲线 $y=f(x)$ 的拐点
\item 点 $x=0$ 不是 $f(x)$ 的极值点，点 $(0,f(0))$ 也不是曲线 $y=f(x)$ 的拐点
\end{choices}
\end{question}
\QuestionBreak

\begin{question}[index=5]
若 $\vect A$ 为三阶实对称阵，$\vect b=(3,-6,-3)^{\mathrm T}$，非齐次线性方程组 $\vect A\vect x=\vect b$ 的通解为
\[
k_1(2,1,0)^{\mathrm T}+k_2(1,0,1)^{\mathrm T}+(-1,2,1)^{\mathrm T},
\]
其中 $k_1,k_2$ 为任意常数，则与矩阵 $\vect A+\vect E$ 合同的为
\begin{choices}[columns=2]
\item $\mat{1&&\\&1&\\&&1}$
\item $\mat{-1&&\\&-1&\\&&-1}$
\item $\mat{-1&&\\&-1&\\&&1}$
\item $\mat{-1&&\\&1&\\&&1}$
\end{choices}
\end{question}
\QuestionBreak

\begin{question}[index=6]
设 $\vect A$ 为 $3$ 阶非零矩阵，$\vect b$ 为 $3$ 维非零列向量，$\vect A^2=\vect O$，则非齐次线性方程组 $\vect A^{\mathrm T}\vect A\vect x=\vect A^{\mathrm T}\vect b$ 线性无关的解向量的个数为
\begin{choices}[columns=4]
\item $1$
\item $2$
\item $3$
\item $4$
\end{choices}
\end{question}
\QuestionBreak

\begin{question}[index=7]
设 $\vect\alpha_1,\vect\alpha_2,\vect\alpha_3$ 为 $n$ 维非零列向量，已知 $\vect\alpha_1-\vect\alpha_2+2\vect\alpha_3=\vect0$，则对任意的 $n$ 维列向量 $\vect\beta$，下列向量组必线性相关的是
\begin{choices}[columns=2]
\item $\vect\beta+\vect\alpha_1,\ 3\vect\beta+\vect\alpha_2,\ 5\vect\beta+\vect\alpha_3$
\item $3\vect\beta+\vect\alpha_1,\ 5\vect\beta+\vect\alpha_2,\ \vect\beta+\vect\alpha_3$
\item $5\vect\beta+\vect\alpha_1,\ \vect\beta+\vect\alpha_2,\ 3\vect\beta+\vect\alpha_3$
\item $\vect\beta+\vect\alpha_1,\ 5\vect\beta+\vect\alpha_2,\ 3\vect\beta+\vect\alpha_3$
\end{choices}
\end{question}
\QuestionBreak

\begin{question}[index=8]
设随机变量 $X$ 与 $Y$ 独立，且 $X\dist E(1)$，$Y\dist U[0,3]$，则 $P\{1<\max(X,Y)\leq2\}=$
\begin{choices}[columns=2]
\item $\dfrac13(1+\eu^{-1}-\eu^{-2})$
\item $\dfrac13(1+\eu^{-1}-2\eu^{-2})$
\item $\dfrac23(1-\eu^{-1}-\eu^{-2})$
\item $\dfrac23(1+\eu^{-1}-2\eu^{-2})$
\end{choices}
\end{question}
\QuestionBreak

\begin{question}[index=9]
设总体 $X\dist E(1)$，$(X_1,X_2,X_3)$ 为来自总体 $X$ 的简单随机样本。记 $\displaystyle T_1=\max_{1\leq i\leq3}X_i,$\quad \allowbreak\mbox{$\displaystyle T_2=\min_{1\leq i\leq3}X_i$}，则 $E(T_1-T_2)=$
\begin{choices}[columns=4]
\item $\dfrac32$
\item $3$
\item $\dfrac43$
\item $1$
\end{choices}
\end{question}
\QuestionBreak

\begin{question}[index=10]
设总体 $X$ 服从参数为 $\lambda$（$\lambda>0$）的泊松分布，$X_1,\dots,X_n$（$n\geq2$）为来自总体 $X$ 的\allowbreak\mbox{简单随机样本}，$\overline X$ 为样本均值，且 $\displaystyle E\left[\sum_{i=1}^n(X_i-\overline X)^2\right]=1-E(\overline X)$，则 $\displaystyle P\left\{\sum_{i=1}^n X_i\geq2\right\}=$
\begin{choices}[columns=4]
\item $\eu^{-1}$
\item $2\eu^{-1}$
\item $1-\eu^{-1}$
\item $1-2\eu^{-1}$
\end{choices}
\end{question}
